% titel: Ballistisches Pendel
% skript: Kap. Stösse, Impulserhaltung, Energieerhaltung

Eine Gewehrkugel bleibt in einem an einer Stange hängenden Klotz stecken
(vollständig inelastischer Stoss); anschliessend schwingt das System als
Pendel hoch. Aus der maximalen Auslenkung wird die Kugelgeschwindigkeit
zurückgerechnet. Kernidee: Man trennt die Probleme "Stoss" und
"Pendelbewegung": der Stoss erhält den Impuls, aber nicht die mechanische
Energie; die Pendelbewegung erhält die Energie, aber nicht den Impuls.

\begin{werkzeugkasten}[Werkzeuge]
\textbf{Inelastischer Stoss} (Kugel $m$, Geschwindigkeit $v$; Ruheklotz $M$;
gemeinsame Geschwindigkeit $u$ nach dem Stoss):
\begin{equation}
  m\,v = (m+M)\,u\;.
\end{equation}

\textbf{Energieerhaltung beim Pendelumlauf} (Höhe $h$, Auslenkung $\varphi$,
Stangenlänge $L$):
\begin{equation}
  \frac12\,(m+M)\,u^2 = (m+M)\,g\,h\;,
  \qquad h = L(1-\cos\varphi)\;.
\end{equation}

\textbf{Überschlag / Seil:} Beim Überschlagen muss im höchsten Punkt die
notwendige Zentripetalbeschleunigung zur Verfügung stehen. Bei einer starren
Stange genügt $u^2/(2)=2gL$; bei einem masselosen Seil muss im höchsten Punkt
zusätzlich $mu_{\mathrm{top}}^2/L > mg$ gelten, sonst schlafft das Seil.

\textbf{Energieverlust durch Deformation:} Die im Stoss "verlorene"
mechanische Energie ist die Differenz der kinetischen Energien vor und nach
dem Stoss.
\end{werkzeugkasten}

\begin{werkzeugkasten}[Vorgehen in allgemeiner Form]
Vom gegebenen $\varphi$ über $h$ auf $u$, dann über die Impulserhaltung auf
$v$ schliessen. Bei Teilaufgaben zur Deformationsenergie die
Energiebilanz unmittelbar vor und nach dem Stoss bilden.
\end{werkzeugkasten}

\begin{achtung}[Typische Fallen]
\begin{itemize}
  \item Die \emph{Gesamtenergie} ist beim inelastischen Stoss nicht erhalten;
       nur der Impuls gilt dort.
  \item Beim elastischen Rückprall des harten Klotzes (Teil d) sind Impuls
       \emph{und} kinetische Energie erhalten.
  \item $h$ nur durch die Pendelgeometrie über $\cos\varphi$ ausdrücken.
\end{itemize}
\end{achtung}